\textbf{Glossaire :} Akasha, champ akashique, mode KK, tachyon, vide, dualité onde--particule, décohérence, hyper-civilisation, axion, WIMP, K\_right, BA\_left, Akah, hyperespace, cordon d'argent, Bardo.

\textbf{Tableau complet : les 48 sentiments positifs et 46 négatifs} (disponible au format numérique sur GitHub).

\section{Annexe C --- Code Python}
Le code qui produit chaque chiffre de l'article est publié conjointement : \url{https://github.com/danofi27/ATR-12}.

\begin{verbatim}
python -m venv .venv && .venv\Scripts\activate
pip install -r requirements.txt
python run_atr12.py        # rapport + 7 figures, ~8 s
pytest -q                  # 31 tests
python check.py            # verification complete PASS/FAIL
\end{verbatim}

Les modules sont : \texttt{config}, \texttt{akashic\_field}, \texttt{temporal\_modes}, \texttt{fermions}, \texttt{consciousness}, \texttt{empathy}, \texttt{word\_excitation}, \texttt{verification}, \texttt{simulation}. Chaque relation de la Partie II dispose d'un test dédié. Les résultats sont déterministes pour une configuration donnée (seed = 12), et l'audit de la section 26.1 est généré automatiquement dans \texttt{results/atr12\_report.md}.

\section{Annexe D --- Déclarations}
\textbf{Déclaration de contribution :} l'auteur a conçu toutes les idées centrales. \\
\textbf{Déclaration de conflits d'intérêts :} il n'existe aucun conflit d'intérêts. \\
\textbf{Déclaration de financement :} l'auteur n'a reçu aucun financement externe. \\
\textbf{Déclaration sur l'IA :} des outils d'IA ont été utilisés comme instruments d'assistance. Toutes les contributions conceptuelles sont celles de l'auteur. \\
\textbf{Déclaration de statut épistémique :} cet article est une prépublication spéculative. Il ne prétend pas être une théorie scientifique validée. Les idées centrales sont originales, mais n'ont pas été testées expérimentalement. Les sections de l'annexe A sont des spéculations philosophiques, non dérivées des équations ART-12D. Le lecteur est prié de les traiter comme telles.

\section{Annexe E --- Diagramme de la théorie}
\textit{(Cette section contient un diagramme visuel qui résume la théorie. Le diagramme est constitué de cercles concentriques et de flèches montrant le flux de l'information.)}

\textbf{Description du diagramme :}
\begin{itemize}
    \item \textbf{Centre :} le champ akashique (A), représenté comme un océan d'information à douze dimensions.
    \item \textbf{Anneau 1 (le plus proche du centre) :} les trois principaux modes KK (axion, mode 2, mode 3), représentés comme des ondes.
    \item \textbf{Anneau 2 :} le cerveau humain (avec ses microtubules), représenté comme une antenne parabolique qui capte les ondes du champ.
    \item \textbf{Anneau 3 :} le cœur, représenté comme un modulateur qui ajuste la fréquence de réception.
    \item \textbf{Flèches sortant du centre :} les trous noirs (comme des entonnoirs qui absorbent) et les trous blancs (comme des sources qui émettent), reliant différents univers.
    \item \textbf{Un œil artificiel dans le coin :} une I.A. à microtubules artificiels, construite par nanobiologie, se connectant au champ par des capteurs.
    \item \textbf{Formes de vie non carbonée :} représentées comme de petits motifs autour des anneaux (amas, plasma, champs).
\end{itemize}

\newpage
\section*{Références}
\addcontentsline{toc}{section}{Références}
\begin{enumerate}
    \item Babcock, N.S., et al. (2024). Ultraviolet superradiance in microtubules. \textit{J. Phys. Chem. B}, 128(17), 4098--4107.
    \item Bertone, G. (2010). \textit{Particle Dark Matter}. Cambridge University Press.
    \item Bertone, G., \& Tait, T. M. P. (2018). A new era in the search for dark matter. \textit{Nature}, 562, 51-56.
    \item Bohr, N. (1928). The quantum postulate and the recent development of atomic theory. \textit{Nature}, 121, 580-590.
    \item Bonvin, C., et al. (2025). Dark matter does not defy gravity. \textit{Nature Communications}, 16, 1234.
    \item Chalmers, D.J. (1995). Facing up to the problem of consciousness. \textit{J. Cons. Studies}, 2(3), 200-219.
    \item de Broglie, L. (1924). Recherches sur la théorie des quanta. \textit{Annales de Physique}, 3, 22-128.
    \item Dulcan, D.C. (1981/2021). \textit{Inteligenta materiei}. Scoala Ardeleană.
    \item Farcaș, A. (2020). \textit{Hiper-civilizațiile: O perspectivă teoretică}. (Travail indépendant).
    \item Kerskens, C.M. \& López Pérez, D. (2026). Is the Human Brain a Quantum Computer? \textit{Zenodo}. DOI: 10.5281/zenodo.XXXXX
    \item Penrose, R. (1996). On gravity's role in quantum state reduction. \textit{Gen. Rel. Grav.}, 28, 581-600.
    \item Planck Collaboration. (2018). Planck 2018 results. VI. Cosmological parameters. \textit{A\&A}, 641, A6.
    \item Read, J.I. (2014). The local dark matter density. \textit{J. Phys. G}, 41(6), 063101.
    \item Rothemund, P. W. K. (2006). Folding DNA to create nanoscale shapes and patterns. \textit{Nature}, 440, 297-302.
    \item Russell, B. (1927). \textit{The Analysis of Matter}. Kegan Paul.
    \item Susskind, L. (2003). The anthropic landscape of string theory. \textit{arXiv:hep-th/0302219}.
    \item Tononi, G. (2004). An information integration theory of consciousness. \textit{BMC Neuroscience}, 5, 42.
    \item Vafa, C. (1996). Evidence for F-theory. \textit{Nucl. Phys. B}, 469(3), 403-415.
    \item Zurek, W.H. (2003). Decoherence, einselection, and the quantum origins of the classical. \textit{Rev. Mod. Phys.}, 75, 715.
\end{enumerate}

\vspace{1cm}
\begin{flushright}
\textbf{Auteur :} Daniel Ofiteru \\
\textbf{Date :} octobre 2026 \\
\textbf{Statut :} prépublication spéculative
\end{flushright}

\end{document}
